\section{データと用語}
\label{sec_data}
対象は，ボケてから収集した6,823,270件のお題である．
各お題には1枚の画像があり，ハッシュ値が同じ画像をまとめると6,609,511枚になる．
以下では，ハッシュ値が同じ画像の1枚を代表の画像と呼び，ラベルの評価と全データの推定は代表の画像を単位として行う．

各お題画像には，形式，属性，媒体の3軸のラベルを付与する．
形式は，回答の作り方を表す1つの値であり，カルタ，穴埋め，お題文，通常の4つがある．
カルタは，ひらがな1文字が入った独立した記号があり，そのひらがなから始まる文を投稿する画像である．
穴埋めは，語の代わりに置かれた空欄があり，空欄に入る語を投稿する画像である．
お題文は，画像中の題に答える画像であり，通常はそれ以外の画像である．
属性は，カルタ記号あり，空欄あり，お題文あり，文字ありの複数選択であり，形式はこの順の優先順位で属性から導く．
媒体は，写真，アニメ，ゲーム，マンガ，テキストグラフィック，絵画，イラストの複数選択である．
本稿で扱うのは，属性のうち文字ありとお題文ありの2項目と，お題文ありから導かれる形式のお題文である．

ラベルの付与は，VLM を教師とし，画像の埋め込みに基づく分類器を生徒とする知識蒸留\cite{ref_kd}で行う．
教師は Qwen3-VL\cite{ref_qwen3vl}であり，vLLM\cite{ref_vllm}で推論する．
生徒は，SigLIP 2\cite{ref_siglip2}の画像埋め込みを入力とする多層パーセプトロン (Multi-Layer Perceptron，MLP) である．
生徒は，教師が約4.5万枚に付けたラベルで学習し，全データに適用する．
以下では，この処理をまとめたプログラムをパッケージと呼び，その版を v5.3.0 のように表す．
また，ラベルの定義を記した文書を定義書と呼び，その版を第1版から第6版のように表す．

正解のラベルは，画像を表示して内容を確認できる AI エージェントである Claude Code が，その時点の定義書に従って目視で付与した．
以下では，このラベルを目視ラベルと呼ぶ．
目視ラベルは全て1名の判断者によるものであり，判断者間の一致度は測っていない．
